\section{伪逆与病态问题}
\label{sec:03-Linear-Algebra/07-SVD-and-Data-Applications/02}

你在拟合一个线性模型，特征表里同时有``身高（厘米）''和``身高（英寸）''两列。两列本该严格成比例，但录入时各自带了一点舍入，所以它们并不完全共线。方程有唯一解，程序不报错，拟合优度也漂亮。换一批数据重跑，两个系数从 \((312,-121)\) 跳到 \((-268,105)\)。预测值几乎没动，系数翻了脸。

系数被什么推着走？沿着``两列同增同减''这个方向调参数，输出几乎不变；换句话说，矩阵把这个输入方向压得极扁。压得越扁，反解时要除的数就越小，数据里千分之一的噪声也能被除成几十倍的系数摆动。SVD 把这个方向的位置指得很准：它就是最小奇异值对应的那个 \(v_i\)。

\subsection{先把``求逆''这个问题问对}

可逆方阵的解是 \(x=A^{-1}b\)。现实里的矩阵经常不方、不满秩，或者因为测量误差而使 \(Ax=b\) 根本无解。这时候继续找逆矩阵，像是要求一张压扁的照片吐出已经丢掉的深度信息。

上一节的分解把该问什么问清楚了。写 \(A=U_r\Sigma_rV_r^{\trans}\)，输出在 \(u_i\) 方向上的分量来自输入在 \(v_i\) 方向上的分量，中间乘了 \(\sigma_i\)。只要 \(\sigma_i>0\)，这一步就能除回去；\(\sigma_i=0\) 的方向已经在正向映射里被抹平，任何公式都变不回来。照这个思路逐方向处理，得到 Moore--Penrose 伪逆

\[
A^{+}=V_r\Sigma_r^{-1}U_r^{\trans}
=\sum_{i=1}^{r}\frac{1}{\sigma_i}\,v_i u_i^{\trans}.
\]

它只反转确实可反转的方向，全程没有假设 \(A\) 可逆，也没有假设它是方阵。

伪逆的价值在于它一次收拾了两种相反的窘境。方程无解时，\(AA^{+}=U_rU_r^{\trans}\) 是到列空间 \(C(A)\) 的正交投影，于是 \(x^{+}=A^{+}b\) 最小化残差 \(\norm{Ax-b}_2\)——这正是\hyperref[sec:03-Linear-Algebra/03-Orthogonality/03]{``最小二乘：为无解方程寻找最佳近似''}给出的解，只是不再要求满列秩。方程解太多时，所有解都写成 \(x^{+}+z\)，其中 \(z\in N(A)\)，而 \(x^{+}\) 落在行空间 \(C(A^{\trans})=N(A)^{\perp}\) 里；两部分正交，添任何非零的 \(z\) 只会把长度撑大，所以 \(x^{+}\) 是范数最小的那个解。

两句话可以合成一句：伪逆在能确定的方向上老实求解，在不能确定的方向上取零。用 \(x_1+x_2=1\) 验一下，无穷多解 \((t,1-t)\) 中它给出 \((1/2,1/2)\)——那个到原点最近的点。

\subsection{有唯一解，不等于解可信}

即使每个奇异值都严格为正，小的那些照样惹祸。考虑

\[
A_\varepsilon=\begin{bmatrix}1&0\\0&\varepsilon\end{bmatrix},
\qquad 0<\varepsilon\ll1 .
\]

第二个方向被压到几乎看不见，反解时却要除以 \(\varepsilon\)。若 \(b\) 的第二分量带着 \(10^{-6}\) 量级的测量误差，而 \(\varepsilon=10^{-8}\)，解的第二分量就会偏出 \(100\)。方程可解，解却无法从有限精度的数据里恢复。

\begin{center}
\begin{tikzpicture}[x=1cm,y=1cm,font=\footnotesize,>=stealth]
  \begin{scope}
    \draw[->] (0,0) -- (0,3.4) node[above] {\(\sigma_i\)};
    \draw[->] (0,0) -- (5.4,0) node[right] {\(i\)};
    \fill[blue!45] (0.5,0) rectangle (1.1,3.00);
    \fill[blue!45] (1.4,0) rectangle (2.0,1.95);
    \fill[blue!45] (2.3,0) rectangle (2.9,1.13);
    \fill[blue!45] (3.2,0) rectangle (3.8,0.30);
    \fill[blue!45] (4.1,0) rectangle (4.7,0.15);
    \node[above] at (0.8,3.02) {\(4\)};
    \node[above] at (4.4,0.17) {\(0.2\)};
    \node at (2.6,-0.7) {正向：奇异值迅速衰减};
  \end{scope}
  \begin{scope}[xshift=7.2cm]
    \draw[->] (0,0) -- (0,3.4) node[above] {\(1/\sigma_i\)};
    \draw[->] (0,0) -- (5.4,0) node[right] {\(i\)};
    \fill[red!45] (0.5,0) rectangle (1.1,0.15);
    \fill[red!45] (1.4,0) rectangle (2.0,0.23);
    \fill[red!45] (2.3,0) rectangle (2.9,0.40);
    \fill[red!45] (3.2,0) rectangle (3.8,1.50);
    \fill[red!45] (4.1,0) rectangle (4.7,3.00);
    \node[above] at (0.8,0.17) {\(0.25\)};
    \node[above] at (4.4,3.02) {\(5\)};
    \node at (2.6,-0.7) {反向：末尾两项主导一切};
  \end{scope}
\end{tikzpicture}

\emph{同一组奇异值 \((4,\,2.6,\,1.5,\,0.4,\,0.2)\)，正向看是一条平缓下滑的曲线，取倒数之后柱子的高低完全颠倒。伪逆把 \(b\) 的噪声按 \(1/\sigma_i\) 分配到各个方向上，末尾那几根柱子决定了解有多不稳。}
\end{center}

要把``放大''说准，得先有矩阵范数。二范数诱导的算子范数 \(\norm{A}_2=\max_{x\ne0}\norm{Ax}_2/\norm{x}_2=\sigma_1\) 量的是最强方向的放大倍数；第二个等号来自把 \(x\) 按右奇异向量展开后 \(\norm{Ax}_2^2=\sum_i\sigma_i^2c_i^2\le\sigma_1^2\)，取 \(x=v_1\) 时相等。Frobenius 范数 \(\norm{A}_F=\sqrt{\sum_{ij}a_{ij}^2}=\sqrt{\sum_i\sigma_i^2}\) 则汇总所有方向的能量，后一等号是因为它在左右乘正交矩阵下不变，可以直接把 \(A\) 换成 \(\Sigma\)。两者回答的不是同一个问题，别因为都叫范数就换着用。

于是可逆方阵的二范数条件数就是奇异值谱的动态范围：

\[
\kappa_2(A)=\norm{A}_2\norm{A^{-1}}_2=\frac{\sigma_1}{\sigma_n},
\]

因为 \(A^{-1}=V\Sigma^{-1}U^{\trans}\) 的奇异值恰是 \(1/\sigma_i\)，最大者为 \(1/\sigma_n\)。它给出的敏感性上界只要两步就能推出来：\(\delta x=A^{-1}\delta b\) 给 \(\norm{\delta x}\le\norm{A^{-1}}_2\norm{\delta b}\)，而 \(\norm{b}=\norm{Ax}\le\norm{A}_2\norm{x}\) 给 \(1/\norm{x}\le\norm{A}_2/\norm{b}\)，两式相乘：

\[
\frac{\norm{\delta x}}{\norm{x}}\le\kappa_2(A)\,\frac{\norm{\delta b}}{\norm{b}} .
\]

这是上界，不是每次都会兑现的预言，它衡量的是问题本身对最坏扰动的敏感程度。算法是否额外引入误差是另一个问题：稳定算法不能让病态问题变得信息充足，只能保证自己不再往下拖（\hyperref[sec:08-Numerical-Analysis/01-Floating-Point-and-Error/03]{``条件数衡量问题本身的敏感性''}把这两笔账分得更细）。

顺带解释一件实践里常被忽略的事。最小二乘可以写成正规方程 \(A^{\trans}Ax=A^{\trans}b\)，但满列秩时 \(A^{\trans}A\) 的特征值是 \(\sigma_i^2\)，于是 \(\kappa_2(A^{\trans}A)=\kappa_2(A)^2\)。原本已经不小的条件数被平方，能用双精度撑住的问题可能就撑不住了。公式在精确算术里成立，不代表它是好算法；实际计算优先走 QR 或 SVD，而不是显式构造 \((A^{\trans}A)^{-1}A^{\trans}\)。

\subsection{截断与岭回归都在做同一件事}

既然麻烦来自除以接近零的数，办法就是别除。最直接的做法是截断：设一个阈值 \(\tau\)，把 \(\sigma_i<\tau\) 的项从伪逆的求和里删掉。这等于宣布那些方向携带的是噪声而非信号，宁可不解，也不要一个被放大一百倍的答案。

更平滑的做法是岭回归。最小化 \(\norm{Ax-b}_2^2+\lambda\norm{x}_2^2\)，把解按奇异向量展开，得到

\[
x_\lambda=\sum_{i}\frac{\sigma_i}{\sigma_i^2+\lambda}\,(u_i^{\trans}b)\,v_i
=\sum_i f_i\cdot\frac{u_i^{\trans}b}{\sigma_i}\,v_i,
\qquad
f_i=\frac{\sigma_i^2}{\sigma_i^2+\lambda}.
\]

把它跟伪逆的 \(\sum_i(u_i^{\trans}b)v_i/\sigma_i\) 摆在一起对比，区别只在那个滤波因子 \(f_i\)。

\begin{center}
\begin{tikzpicture}[x=1.6cm,y=1.2cm,font=\footnotesize,>=stealth]
  \draw[->] (0,0) -- (3.3,0) node[right] {\(\sigma\)};
  \draw[->] (0,0) -- (0,2.5) node[above] {\(f=\sigma^2/(\sigma^2+\lambda)\)};
  \draw[dashed,black!50] (0,2) -- (3.2,2);
  \node[left] at (0,2) {\(1\)};
  \draw[blue!70!black,thick] plot[domain=0:3.2,samples=90] (\x,{2*(\x*\x)/(\x*\x+0.25)});
  \draw[dotted,black!60] (0.5,0) -- (0.5,1);
  \draw[dotted,black!60] (0,1) -- (0.5,1);
  \node[below] at (0.5,-0.05) {\(\sqrt\lambda\)};
  \node[left] at (0,1) {\(1/2\)};
  \node[align=left,black!70] at (2.15,0.85) {大奇异值方向\\几乎原样保留};
\end{tikzpicture}

\emph{滤波因子在 \(\sigma\gg\sqrt\lambda\) 处趋于 1，方向被原样保留；在 \(\sigma\ll\sqrt\lambda\) 处约等于 \(\sigma^2/\lambda\)，那一项的系数 \(f_i/\sigma_i\approx\sigma_i/\lambda\) 随 \(\sigma_i\) 一起趋零，而不是随 \(1/\sigma_i\) 爆炸。截断是这条曲线的硬开关版本。}
\end{center}

换句话说，\(\lambda\) 给奇异值垫了一个底，谁也不会被除到无穷去。后果写在脸上：所有方向都被按住了一点，包括那些本该原样保留的，于是解有了偏差。稳定性与偏差的这笔交换在\hyperref[sec:09-Convex-Optimization/07-Applications/01]{``光滑经验风险与 \(L_2\) 正则化''}里会从优化的角度重讲一遍。机器学习里的过拟合与这里的病态是同一种病的两种说法：模型沿着数据几乎没有约束的方向自由摆动，正则化的作用就是把那些方向按住。

最后收几条容易踩空的地方。伪逆不是把矩阵元素逐个取倒数，也不总等于 \((A^{\trans}A)^{-1}A^{\trans}\)，后者要求满列秩。\(A^{+}A\) 投影到行空间，\(AA^{+}\) 投影到列空间，两者尺寸都不同。条件数大只说明问题敏感，不证明你这一次算错了；条件数小也不豁免建模误差。最需要警惕的是秩的断崖：\(\sigma=0\) 与 \(\sigma=10^{-14}\) 在数学上分属两个世界，伪逆会因此突变，而浮点数分不清它们——阈值 \(\tau\) 怎么定，因此是个需要交代的建模决定，而不是可以留给默认参数的细节。

截断把小奇异值方向整个扔掉，理由到这里还只是``它们不可信''。下一节换个问法：如果只保留前 \(k\) 个方向，损失能否算准，以及在所有秩不超过 \(k\) 的矩阵里，这个截断是不是最好的那一个。
